\subsection{Two lecture blocks and one central idea}

The chapter is not a catalogue of matrix notation. Its central question is: \emph{how can one rectangular array organise several related quantities and operations at once?}

\begin{corebox}[First 45 minutes: reading and combining matrices]
The essential material is the shape of a matrix, the meaning of an entry, addition and scaling, and the row--column mechanism of matrix multiplication. The main goal is to read every operation structurally rather than to memorise a formula.
\end{corebox}

\begin{corebox}[Second 45 minutes: structure that will be reused]
Identity, diagonal, triangular, symmetric, and transpose patterns show how visible form simplifies behaviour. Elementary row operations are introduced as reversible changes that will later acquire three different meanings: determinant bookkeeping, construction of an inverse, and replacement of a linear system by an equivalent one.
\end{corebox}

The later chapters return to row operations deliberately. Each return adds a new interpretation; it is not a repetition of an isolated algorithm.
